\documentclass[10pt,letterpaper,oneside]{article}
\usepackage[T1]{fontenc}
\usepackage[utf8]{inputenc}
\usepackage{lmodern}
\usepackage{geometry}
\geometry{margin=0.84in,headheight=21pt,footskip=28pt}
\usepackage{microtype}
\usepackage{xcolor}
\definecolor{ink}{HTML}{15253D}
\definecolor{accent}{HTML}{086A84}
\definecolor{pale}{HTML}{EEF5F8}
\definecolor{muted}{HTML}{4E6474}
\definecolor{grayline}{HTML}{D7E2E8}
\usepackage{graphicx}
\usepackage{fancyhdr}
\usepackage{titlesec}
\usepackage{enumitem}
\usepackage{array,tabularx,longtable,booktabs}
\usepackage[most]{tcolorbox}
\usepackage{xurl}
\usepackage{hyperref}
\usepackage{listings}
\usepackage{amsmath,amssymb}
\usepackage{tikz}
\usetikzlibrary{arrows.meta,positioning,shapes.geometric,calc}
\hypersetup{colorlinks=true,linkcolor=accent,urlcolor=accent,pdftitle={Governance, Risk and Resilience},pdfauthor={Independent technical engineering study companion}}
\renewcommand{\familydefault}{\sfdefault}
\linespread{1.07}
\setlength{\parindent}{0pt}
\setlength{\parskip}{5.5pt}
\emergencystretch=2em
\setlist[itemize]{leftmargin=1.5em,itemsep=2.5pt,topsep=3pt}
\setlist[enumerate]{leftmargin=1.7em,itemsep=3pt,topsep=3pt}
\titleformat{\section}{\Large\bfseries\color{ink}}{\thesection}{0.65em}{}[\vspace{-2pt}{\color{accent}\titlerule}]
\titleformat{\subsection}{\large\bfseries\color{accent}}{\thesubsection}{0.6em}{}
\titleformat{\subsubsection}{\bfseries\color{ink}}{\thesubsubsection}{0.6em}{}
\titlespacing*{\section}{0pt}{16pt}{9pt}
\titlespacing*{\subsection}{0pt}{11pt}{6pt}
\pagestyle{fancy}\fancyhf{}\fancyhead[L]{\footnotesize\color{muted} COMPUTER \& INFORMATION SECURITY ENGINEERING}\fancyhead[R]{\footnotesize\color{accent} SUITE / 3}\fancyfoot[L]{\footnotesize\color{muted} Independent study companion \quad|\quad 2026}\fancyfoot[R]{\footnotesize\color{muted}\thepage}\renewcommand{\headrulewidth}{0.4pt}\renewcommand{\headrule}{\hbox to\headwidth{\color{grayline}\leaders\hrule height \headrulewidth\hfill}}
\newtcolorbox{sourcebox}{enhanced,breakable,boxrule=0pt,leftrule=3pt,colframe=accent,colback=pale,arc=0pt,left=9pt,right=7pt,top=7pt,bottom=7pt}
\newtcolorbox{warningbox}{enhanced,breakable,boxrule=0.5pt,colframe=grayline,colback=white,arc=1pt,left=9pt,right=9pt,top=7pt,bottom=7pt}
\newcommand{\handbookref}[1]{\textsf{\color{accent}\textbf{[H:#1]}}}
\newcommand{\modernext}{\textsf{\color{muted}\textbf{[E]}}}
\newcommand{\designed}{\textsf{\color{accent}\textbf{[D]}}}
\newcommand{\deliverable}[1]{\textbf{Deliverable:} #1}
\lstdefinestyle{lab}{basicstyle=\ttfamily\scriptsize,keywordstyle=\color{accent}\bfseries,commentstyle=\color{muted},numbers=left,numberstyle=\tiny\color{muted},numbersep=7pt,breaklines=true,showstringspaces=false,frame=single,rulecolor=\color{grayline},backgroundcolor=\color{pale},xleftmargin=1em,framexleftmargin=0.6em,tabsize=2,columns=fullflexible}
\lstset{style=lab}
\renewcommand{\arraystretch}{1.21}
\newcommand{\makereportcover}[4]{%
\thispagestyle{empty}\vspace*{1.15in}
{\small\bfseries\color{accent} COMPUTER AND INFORMATION SECURITY ENGINEERING}\par\vspace{13pt}
{\Huge\bfseries\color{ink} #1\par}\vspace{15pt}
{\Large\color{muted} #2\par}\vspace{18pt}
\color{accent}\rule{0.94\textwidth}{2.1pt}\par\vspace{20pt}
{\normalsize\bfseries Document #3 of 8 \quad\textbullet\quad 16 October 2026\par}
\vspace{0.35in}
\begin{sourcebox}\textbf{Orientation.} #4\end{sourcebox}
\vfill
{\small\color{muted} Independent, original engineering study and practice materials grounded in\par
\emph{Computer and Information Security Handbook}, 4th ed., edited by John R. Vacca,\par
Elsevier/Morgan Kaufmann, 2025, Volumes I--II. Not a reproduction or official standard.}
\newpage\setcounter{page}{1}\tableofcontents\newpage
}
\begin{document}
\makereportcover{Governance, Risk and Resilience}{Policies, control systems, metrics, audit, continuity and recovery}{3}{An operational management reference with templates for risk registers, policy approval, evidence and business-impact analysis; modern frameworks are cross-references, not substituted for the handbook.}

\section{Governance as a managed system}
The handbook distinguishes executive alignment, security management systems, risk-oriented security management, policies, education, metrics, certification and disaster recovery. \handbookref{1,2,24--27,33--35,37--40,81} An effective security program needs decisions, owners, escalation paths, evidence and feedback---not merely a set of policies.
\begin{sourcebox}\textbf{Recommended management cadence} \designed{}: strategy and risk appetite (annual or after major change); risk review (monthly/quarterly according to exposure); exceptions (before expiry); control testing (risk-based); disaster recovery exercise (scheduled); incidents and material changes (event-driven). Calendars are program decisions, not universal mandates.\end{sourcebox}

\section{Governance and operating model}
\subsection{Accountabilities}
\par\medskip\noindent\begin{tabularx}{\textwidth}{@{}p{0.25\textwidth}X@{}}\toprule
Accountability & Decision and evidence \\ \midrule
Governing body & Approves risk appetite; receives enterprise risk and resilience reports.\\
Business asset owner & Classifies asset; accepts or rejects residual business risk.\\
Security program leader & Maintains strategy, standards, program metrics and challenge process.\\
Control operator & Operates safeguards and provides dated evidence.\\
Independent assurance & Tests whether design and operation conform to stated criteria.\\
Incident commander & Coordinates response within delegated authority.\\
Privacy/legal & Identifies applicable processing, retention and notification obligations.\\
\bottomrule\end{tabularx}\par\medskip
The handbook's people--process--technology framing and standards integration inform this split. The roles above are a designed RACI model, not claimed as prescribed by a chapter. \handbookref{2,25,27} \designed{}

\subsection{Policy hierarchy}\designed{}
\begin{enumerate}
\item \textbf{Policy:} leadership intent and non-negotiable obligations.
\item \textbf{Standard:} measurable minimum required baseline (e.g., admin MFA).
\item \textbf{Procedure:} operational sequence for performing a control.
\item \textbf{Guideline:} recommended implementation patterns and exceptions.
\item \textbf{Record:} evidence of an approval, test, configuration or decision.
\end{enumerate}
Every statement requiring compliance must indicate its source of authority, owner, applicable population, effective date, review date and deviation route. The handbook discusses creation and approval of security policies and standards implementation. \handbookref{39,40}

\section{Risk assessment and treatment}
\subsection{Scenario-based register}\handbookref{35}
\par\medskip\noindent\begin{longtable}{@{}p{0.14\textwidth}p{0.28\textwidth}p{0.15\textwidth}p{0.26\textwidth}@{}}\toprule
ID & Risk scenario & Business outcome & Treatment and proof \\ \midrule\endhead
R-01 & Internet-facing service has exploitable authorization defect. & Unauthorized disclosure. & Fix in code; deny cross-tenant test, reviewed deployment.\\
R-02 & Admin identities lack phishing-resistant MFA. & Account takeover. & Enforce policy; authentication method report.\\
R-03 & Production backup uses same administrative credentials. & Irrecoverable loss. & Isolate access; verify destructive-access test and restore.\\
R-04 & Safety network accepts unsanctioned commands. & Equipment damage or injury. & Network isolation, command validation, fail-safe simulation.\\
\bottomrule\end{longtable}\par\medskip

For every risk record add asset identifier, threat actor or fault, observed weakness, impacted stakeholders, inherent exposure, controls, residual assessment, treatment owner, target date, validation owner and acceptance authority. Avoid treating a color-coded risk score as an objective measurement. \designed{}
\subsection{Simple quantitative illustrations}\designed{}
If a loss event with consequence \(L\) has estimated annualized frequency \(\lambda\), a simplified expected annual loss is
\[\mathrm{EAL}=\lambda L.\]
Example: a \$250,000 modeled loss with \(\lambda=0.08\) gives \$20,000 expected annual loss. This is a \emph{scenario assumption}, not a forecast or justification for risk acceptance. Record uncertainty ranges, nonfinancial impacts, tail losses, detection gaps and dependencies. A low average loss can coexist with intolerable safety or legal exposure.

\subsection{Risk treatment decision record}\designed{}
\begin{warningbox}\textbf{RT-XXX:} scenario, owner, date and scope; alternatives \emph{avoid, mitigate, transfer, accept}; selected treatment and cost; expected change to likelihood/impact; test of effectiveness; residual risk; expiry/re-evaluation trigger; signature of delegated accepting authority. Acceptance is not an engineering waiver: retain the finding and its audit trail.\end{warningbox}

\section{Framework alignment without false equivalence}
The handbook explicitly references ISO/IEC 27002:2022, standards integration and NIST CSF-era practices. \handbookref{2,24,35,39} \modernext{} NIST CSF 2.0 organizes outcomes into Govern, Identify, Protect, Detect, Respond and Recover. This is an outcomes vocabulary, not a one-to-one replacement for a system's control standard.
\par\medskip\noindent\begin{tabularx}{\textwidth}{@{}p{0.16\textwidth}X X@{}}\toprule
CSF function & Program question & Example evidence \\ \midrule
Govern & Who decides and owns cyber risk? & Risk appetite, policy approvals\\
Identify & What assets and exposure matter? & System inventory, BIA, dependencies\\
Protect & Which safeguards limit harm? & Approved baseline, access reviews\\
Detect & Can important deviations be recognized? & Coverage report and detection test\\
Respond & Can responders contain and communicate? & Tabletop record, playbooks\\
Recover & Can services be restored? & Restore test and improvement log\\
\bottomrule\end{tabularx}\par\medskip

\section{Control validation and meaningful metrics}
Security metrics must reveal decisions and changes rather than supply attractive but unhelpful dashboards; the handbook's metrics chapter is explicitly concerned with their selection and usefulness. \handbookref{33}
\par\medskip\noindent\begin{tabularx}{\textwidth}{@{}p{0.32\textwidth}X@{}}\toprule
Metric & Definition and interpretation \\ \midrule
Asset-owner completeness & Assets with an accountable owner / in-scope assets.\\
Critical finding SLA compliance & Critical findings resolved within policy interval / critical findings due.\\
Detection coverage & Modeled priority attack paths with tested detection / all modeled priority paths.\\
Recovery success & Restore scenarios completed within stated RTO/RPO / those exercised.\\
Exception debt & Accepted exceptions past expiry, with impact-weighted review.\\
Mean time to contain & Time from confirmed incident to verified containment, with segmentation by severity.\\
\bottomrule\end{tabularx}\par\medskip
\designed{} Pair rates with denominators, selection criteria, blind spots and raw evidence. A falling finding count may indicate a shrinking scan scope, not improving security. Distinguish leading indicators (control deployment) from lagging indicators (real losses and incidents).

\section{Disaster recovery and business continuity}
\handbookref{9,37,38,90} The source emphasizes risk measurement, business impact analysis, recovery planning, training, tests and communications. Engineering converts these into service-specific recovery objectives.
\begin{description}[leftmargin=1.6em,style=nextline]
\item[RTO] Target elapsed time to restore an acceptable level of service after disruption.
\item[RPO] Maximum tolerable data-loss interval, stated in terms of recoverable data state.
\item[MTD / MTPD] Maximum tolerable duration of disruption for a defined mission or process.
\end{description}
\subsection{Example recovery dependency table}\designed{}
\par\medskip\noindent\begin{tabularx}{\textwidth}{@{}p{0.19\textwidth}p{0.17\textwidth}p{0.17\textwidth}X@{}}\toprule
Service & Target RTO & Target RPO & Restoration dependency \\ \midrule
Identity and secrets & 2 hours & 15 minutes & Isolated restore identity and protected keys\\
Customer portal & 4 hours & 1 hour & DNS, gateway, auth, data store\\
Analytics reporting & 48 hours & 24 hours & Event stream and restored warehouse\\
\bottomrule\end{tabularx}\par\medskip
These are \emph{illustrative}, not generally appropriate thresholds. The business owner must approve objectives consistent with impact and technical feasibility.
\subsection{Recovery exercise protocol}
\begin{enumerate}
\item Approve scope, recovery scenario, communications and safety guardrails.
\item Capture baseline state and backup timestamps.
\item Simulate loss with a nonproduction environment or a tightly controlled production window.
\item Recover credentials, infrastructure, application and data in dependency order.
\item Measure outage, data loss, correctness, performance and security-control availability.
\item Record deviations, remediation owners and repeat test criteria.
\end{enumerate}

\section{Audit, compliance and improvement}
\handbookref{39,40,81} An audit tests controls against \emph{defined criteria}. Do not label an organization ``compliant'' merely because a scanner passed. Preserve policy version, applicable control, sample population, tester independence, sample dates, exceptions and auditor conclusion.
\begin{sourcebox}\textbf{Evidence packet template} \designed{}: control ID; control objective; owner; implementation narrative; population and boundaries; test procedure; test date; results; links to immutable records; failures; remediation ticket; approval.\end{sourcebox}

\section{Practice: quarterly risk committee}
Use an imagined cloud-hosted healthcare service. Prepare a risk register with eight scenarios, assign ownership, justify three investment decisions, create a five-metric dashboard and plan a recovery test. Grading: accountable decisions (25), traceable criteria (25), realistic dependencies (20), evidence design (20), residual-risk communication (10). \designed{}

\section*{Handbook chapter alignment}\addcontentsline{toc}{section}{Handbook chapter alignment}
These chapters ground the domain scope. The document offers original synthesis and technical procedures, not a chapter-by-chapter reproduction.
\par\medskip\noindent\begin{longtable}{@{}p{0.085\textwidth}p{0.68\textwidth}p{0.16\textwidth}@{}}\toprule Chapter & Handbook topic & Combined PDF p. \\ \midrule\endhead1 & Information Security in the Modern Enterprise & 39 \\
2 & Building a Secure Organization & 49 \\
24 & Information Security Essentials for IT Managers: Protecting Mission-Critical Systems & 459 \\
25 & Security Management Systems & 469 \\
26 & Policy-Driven System Management & 477 \\
27 & Information Technology Security Management & 511 \\
33 & Security Metrics: An Introduction and Literature Review & 595 \\
34 & Security Education, Training, and Awareness & 609 \\
35 & Risk Management & 619 \\
36 & Insider Threats & 637 \\
37 & Disaster Recovery & 651 \\
38 & Disaster Recovery Plans for Small and Medium Business (SMB) & 661 \\
39 & Security Certification and Standards Implementation & 671 \\
40 & Security Policies and Plans Development & 679 \\
81 & Assessments and Audits & 1369 \\
90 & Disaster Management Cycle, Smart Technologies, Urban Form, Traditional Developments, and Sprawl & 1499 \\
95 & Online e-Reputation Management Services & 1595 \\
97 & Data Loss Protection & 1629 \\
\bottomrule\end{longtable}\par\medskip

\section*{Source and reference register}
\addcontentsline{toc}{section}{Source and reference register}
\begin{description}[leftmargin=1.6em,style=nextline]
\item[H] John R. Vacca (ed.), \emph{Computer and Information Security Handbook}, 4th ed., Volumes I--II, Elsevier/Morgan Kaufmann, 2025. Chapter identifiers in this document point to this book.
\item[N1] NIST, \emph{Cybersecurity Framework 2.0} (2024), \url{https://www.nist.gov/cyberframework}.
\item[N2] NIST SP 800-218, \emph{Secure Software Development Framework v1.1} (2022), \url{https://csrc.nist.gov/pubs/sp/800/218/final}. Later drafts are not treated as final requirements.
\item[N3] NIST SP 800-61 Rev. 3, \emph{Incident Response Recommendations and Considerations for Cybersecurity Risk Management} (2025), \url{https://csrc.nist.gov/pubs/sp/800/61/r3/final}.
\item[N4] NIST SP 800-63-4, \emph{Digital Identity Guidelines} (2025), \url{https://csrc.nist.gov/pubs/sp/800/63/4/final}.
\item[N5] NIST, FIPS 203/204/205 (2024), \url{https://csrc.nist.gov/Projects/Post-Quantum-Cryptography/publications}.
\item[O1] OWASP, \emph{API Security Top 10} (2023), \url{https://api-security.owasp.org/editions/2023/en/0x11-t10/}.
\item[C1] CISA, \emph{Known Exploited Vulnerabilities Catalog}, \url{https://www.cisa.gov/known-exploited-vulnerabilities-catalog}.
\item[M1] MITRE ATT\&CK Enterprise, \url{https://attack.mitre.org/matrices/enterprise/}.
\item[G1] GitHub, \emph{REST API endpoints for code scanning}, \url{https://docs.github.com/en/rest/code-scanning}.
\end{description}
\begin{warningbox}\textbf{Interpretation rule.} \handbookref{number} identifies a conceptual anchor in the supplied handbook, not a verbatim quotation or an official requirement. \modernext{} identifies a contemporary extension informed by external references or current engineering practice. \designed{} identifies an independently developed example, template, exercise, metric, or operational procedure. Requirements are binding only if separately adopted through applicable law, contract, policy, or the cited standards; this suite does not assert otherwise.\end{warningbox}

\end{document}
